\documentclass[12pt,a4paper]{ctexart}

\usepackage[margin=2.4cm]{geometry}
\usepackage{enumitem}

\setlist[itemize]{leftmargin=2.2em,itemsep=0.35em,topsep=0.35em}
\setlist[enumerate]{leftmargin=2.4em,itemsep=0.25em,topsep=0.25em}

\title{《习概选择题》涉及会议知识点整理}
\author{}
\date{}

\begin{document}
\section*{一、党的重要代表大会}

\subsection*{党的十二大}

\begin{itemize}
\item 首次把“小康”作为经济建设的总奋斗目标。
\end{itemize}

\subsection*{党的十五大}

\begin{itemize}
\item 明确提出形成中国特色社会主义法律体系的重大任务。
\end{itemize}

\subsection*{党的十六大}

\begin{itemize}
\item 针对小康社会低水平、不全面、发展很不平衡的实际，提出全面建设小康社会目标。
\end{itemize}

\subsection*{党的十八大以来（题库中的阶段性表述）}

\begin{itemize}
\item 以中国式现代化全面推进中华民族伟大复兴。
\item 中国特色社会主义进入新时代，成为我国发展新的历史方位。
\item 在实现第一个百年奋斗目标的基础上，开启向第二个百年奋斗目标进军的新征程。
\item 坚持以人民为中心的发展思想，在发展中保障和改善民生，加强和创新社会治理。
\end{itemize}

\subsection*{党的十九大}

\begin{itemize}
\item 通过《关于〈中国共产党章程（修正案）〉的决议》，将习近平新时代中国特色社会主义思想确立为党必须长期坚持的指导思想，并写入党章。
\item 概括以“十四个坚持”为主要内容的基本方略。
\item 提出到2035年各方面制度更加完善，国家治理体系和治理能力现代化基本实现。
\item 将“增强绿水青山就是金山银山的意识”等内容写入党章。
\item 提出到本世纪中叶把人民军队全面建成世界一流军队。
\item 在报告中进一步强调“打铁必须自身硬”。
\end{itemize}

\subsection*{党的二十大}

\begin{itemize}
\item 提出“六个必须坚持”，即：
\begin{enumerate}
\item 必须坚持人民至上；
\item 必须坚持自信自立；
\item 必须坚持守正创新；
\item 必须坚持问题导向；
\item 必须坚持系统观念；
\item 必须坚持胸怀天下。
\end{enumerate}
\item “六个必须坚持”是习近平新时代中国特色社会主义思想的世界观、方法论和贯穿其中的立场观点方法的重要体现，彰显其理论品格和鲜明特征。
\item 围绕全面建成社会主义现代化强国、以中国式现代化全面推进中华民族伟大复兴的中心任务，对全面深化改革开放作出新部署。
\item 强调统筹发展和安全，并将其写入党章。
\end{itemize}

\section*{二、党的中央全会}

\subsection*{党的十八届三中全会（2013年11月）}

\begin{itemize}
\item 审议通过《中共中央关于全面深化改革若干重大问题的决定》，对全面深化改革进行总体擘画。
\item 确定全面深化改革总目标：完善和发展中国特色社会主义制度，推进国家治理体系和治理能力现代化。
\item 首次提出“推进法治中国建设”的目标，并作出重大任务部署，更加突出法治国家、法治政府、法治社会一体建设。
\item 提出创新社会治理体制，推动由“社会管理”向“社会治理”转变。
\end{itemize}

\subsection*{党的十八届四中全会}

\begin{itemize}
\item 审议通过《中共中央关于全面推进依法治国若干重大问题的决定》。
\item 对全面推进依法治国作出总体部署，规划法治中国建设的总蓝图、路线图和施工图。
\end{itemize}

\subsection*{党的十九届五中全会}

\begin{itemize}
\item 进一步充实国防和军队现代化的目标任务和发展步骤。
\item 明确提出确保2027年实现建军一百年奋斗目标。
\end{itemize}

\subsection*{党的十九届六中全会}

\begin{itemize}
\item 提出“两个确立”：确立习近平同志党中央的核心、全党的核心地位；确立习近平新时代中国特色社会主义思想的指导地位。
\item 审议通过《中共中央关于党的百年奋斗重大成就和历史经验的决议》。
\item “十个明确”集中体现习近平新时代中国特色社会主义思想的主要观点和基本精神，是其主体内容和“四梁八柱”。
\item 将“发展全过程人民民主”作为习近平新时代中国特色社会主义思想的重要内容。
\end{itemize}

\subsection*{党的二十届三中全会}

\begin{itemize}
\item 于2024年7月18日在北京闭幕。
\item 重点研究部署进一步全面深化改革、推进中国式现代化问题。
\item 审议通过《中共中央关于进一步全面深化改革、推进中国式现代化的决定》。
\item 明确改革开放是中国式现代化的必由之路，党的领导是中国式现代化的根本保证。
\item 强调进一步全面深化改革，要进一步解放和发展社会生产力、激发和增强社会活力，以促进社会公平正义、增进人民福祉为出发点和落脚点。
\end{itemize}

\paragraph{2035年目标}

\begin{itemize}
\item 全面建成高水平社会主义市场经济体制；
\item 中国特色社会主义制度更加完善；
\item 基本实现国家治理体系和治理能力现代化；
\item 基本实现社会主义现代化；
\item 为到本世纪中叶全面建成社会主义现代化强国奠定坚实基础。
\end{itemize}

\paragraph{经济体制与高质量发展}

\begin{itemize}
\item 高质量发展是全面建设社会主义现代化国家的首要任务。
\item 健全促进实体经济和数字经济深度融合制度。
\item 完善强农惠农富农支持制度，坚持农业农村优先发展，完善乡村振兴投入机制。
\item 推进高水平对外开放，建设现代化经济体系。
\item 加快构建新发展格局，推动高质量发展。
\item 坚持和落实“两个毫不动摇”：毫不动摇巩固和发展公有制经济；毫不动摇鼓励、支持、引导非公有制经济发展。
\end{itemize}

\paragraph{教育、科技、人才与全过程人民民主}

\begin{itemize}
\item 教育、科技、人才是中国式现代化的基础性、战略性支撑。
\item 丰富各层级民主形式。
\item 健全全过程人民民主制度体系：坚持党的领导、人民当家作主、依法治国有机统一；推动人民当家作主制度更加健全；促进协商民主广泛多层制度化发展；促进中国特色社会主义法治体系更加完善；推动社会主义法治国家建设达到更高水平。
\end{itemize}

\paragraph{文化体制机制改革}

\begin{itemize}
\item 构建更有效力的国际传播体系，推进国际传播格局重构。
\item 优化文化服务和文化产品供给机制。
\item 完善公共文化服务体系，建立优质文化资源直达基层机制。
\item 健全社会力量参与公共文化服务机制，推进公共文化设施所有权和使用权分置改革。
\item 深化文化领域国资国企改革。
\end{itemize}

\paragraph{民生、收入分配与社会保障}

\begin{itemize}
\item 健全基本公共服务制度体系，加强普惠性、基础性、兜底性民生建设。
\item 构建初次分配、再分配、第三次分配协调配套的制度体系，提高居民收入在国民收入分配中的比重和劳动报酬在初次分配中的比重。
\item 完善税收、社会保障、转移支付等再分配调节机制，支持发展公益慈善事业。
\item 支持和规范发展新就业形态。
\item 按照自愿、弹性原则，稳妥有序推进渐进式延迟法定退休年龄改革。
\item 促进医疗、医保、医药协同发展和治理。
\item 完善基本养老保险全国统筹制度，健全全国统一的社会保障公共服务平台。
\item 健全社会保障基金保值增值和安全监管体系，健全基本养老、基本医疗保险筹资和待遇合理调整机制，逐步提高城乡居民基本养老保险基础养老金。
\end{itemize}

\paragraph{生态文明体制改革}

\begin{itemize}
\item 实施分区域、差异化、精准管控的生态环境管理制度。
\item 健全自然资源资产产权制度和管理制度体系，完善全民所有自然资源资产所有权委托代理机制。
\item 建立生态环境保护、自然资源保护利用和资产保值增值等责任考核监督制度。
\item 完善国家生态安全工作协调机制。
\end{itemize}

\section*{三、其他国内重要会议和大会}

\subsection*{2017年12月中央经济工作会议}

\begin{itemize}
\item 首次对习近平经济思想作出系统概括。
\end{itemize}

\subsection*{十三届全国人大一次会议（2018年3月）}

\begin{itemize}
\item 通过宪法修正案，郑重将习近平新时代中国特色社会主义思想载入宪法。
\end{itemize}

\subsection*{中央国家安全委员会第一次会议（2014年4月15日）}

\begin{itemize}
\item 首次提出总体国家安全观，强调坚持总体国家安全观、走中国特色国家安全道路。
\end{itemize}

\subsection*{中央军委扩大会议（2012年12月）}

\begin{itemize}
\item 提出走中国特色强军之路，为建设一支听党指挥、能打胜仗、作风优良的人民军队而奋斗。
\end{itemize}

\subsection*{庆祝中国共产党成立100周年大会（2021年7月1日）}

\begin{itemize}
\item 庄严宣告实现第一个百年奋斗目标，在中华大地上全面建成小康社会，历史性地解决绝对贫困问题。
\item 宣示向全面建成社会主义现代化强国的第二个百年奋斗目标迈进。
\end{itemize}

\subsection*{十三届全国人大常委会第三十八次会议（2022年12月）}

\begin{itemize}
\item 通过《关于〈中华人民共和国香港特别行政区维护国家安全法〉第十四条和第四十七条的解释》。
\item 及时妥善解决香港国安法实施中的实际问题，确保香港国安法正确、有效实施。
\end{itemize}

\section*{四、题库涉及的国际大会}

\subsection*{联合国《生物多样性公约》第十五次缔约方大会（COP15）}

\begin{itemize}
\item 中国成功承办会议，推动达成“昆明--蒙特利尔全球生物多样性框架”等一揽子具有里程碑意义的成果。
\item 为全球生物多样性治理确定目标、明确路径。
\end{itemize}

\subsection*{第七十届联合国大会一般性辩论（2015年9月）}

\begin{itemize}
\item 提出和平、发展、公平、正义、民主、自由是全人类共同价值，也是联合国的崇高目标。
\end{itemize}

\subsection*{第七十六届联合国大会（2021年9月）}

\begin{itemize}
\item 提出全球发展倡议。
\end{itemize}

\end{document}
